\begin{abstract}
IT support teams must decide what to do with an incident before they know its cause: investigate further, apply a fix, or pass the incident to someone better placed to resolve it. Reinforcement learning (RL) can treat this as a sequence of decisions, but the published work is spread across IT operations, security operations and safe RL. This review asks how RL and related sequential methods have been applied to IT incident handling, which of the three decisions their agents make, how operational rules are enforced, and how the agents are evaluated.

The review follows the PRISMA 2020 guideline. Searches of Scopus, IEEE Xplore and the ACM Digital Library on 23 October 2026 returned \NIdentified{} records. After the removal of \NDuplicates{} duplicates and \NAuto{} records by automated filtering, \NScreened{} records were screened on title and abstract, \NSought{} reports were sought, \NAssessed{} were assessed in full, and \NIncluded{} studies published between \NYearMin{} and 2026 were included. Their quality was appraised with an eight-item checklist, and the evidence was synthesised narratively around five review questions.

There are four main findings. First, only \NBelief{} studies represent the hidden state of an incident with an explicit belief, and most treat it as fully observed. Second, \NInfo{} studies give the agent an explicit information-gathering action, but only \NEsc{} include escalation to a person, and only one combines gathering, a response action and escalation, with the escalation fixed by design rather than learned. Third, rule enforcement is proposed more often than it is measured: \NShield{} studies use a hard filter or shield, and \NBlocked{} report what happens to a blocked or uncertain decision. Fourth, decision quality is measured in \NDQ{} studies but escalation behaviour in \NEscM{} and constraint compliance in \NComp{}, and robustness to changing conditions is tested in \NRobust{}. The evidence is recent, with \NYearRecent{} studies published in 2025 or 2026, and uneven in quality: \NQLow{} studies meet at most two of the eight checklist items. The review concludes that a single policy which learns to investigate, act or escalate for IT service incidents under explicit rules has not been reported, and it sets out the evaluation that such work should include.
\end{abstract}
